\section{La suma (o adición)}

Ahora que se han introducido los números naturales, se puede estudiar una de las operaciones más elementales que realizamos con ellos: la adición, comúnmente llamada suma.

La adición toma dos números naturales y les asigna otro número natural. Esta idea puede expresarse simbólicamente como
\[
+\colon\mathbb{N}\times\mathbb{N}\to\mathbb{N}.
\]
No se preocupe si esta notación parece extraña a primera vista. Por ahora solo queremos expresar, de una manera compacta y precisa, que la adición recibe dos números naturales y produce un número natural. A continuación analizaremos cada parte de esta escritura.

\subsection{Los números naturales}

El símbolo \(\mathbb{N}\) representa el conjunto de los números naturales, es decir, los números que utilizamos principalmente para contar:
\[
1,2,3,4,5,\dots
\]
Dependiendo de la convención utilizada, el cero también puede considerarse un número natural. En este libro adoptaremos la convención
\[
\mathbb{N}=\{0,1,2,3,\dots\}.
\]
El conjunto de los números naturales es infinito: no existe un último número natural. Por grande que sea el número que elija, siempre habrá números naturales mayores que él.

\begin{note}
En matemáticas es importante prestar atención a palabras como \emph{siempre}, \emph{nunca}, \emph{todos} o \emph{ninguno}. Estas palabras expresan afirmaciones generales y, por lo tanto, deben cumplirse en todos los casos comprendidos por la afirmación.

Por ejemplo, decir que siempre existe un número natural mayor que otro significa precisamente que esto debe ser cierto sin importar qué número natural se haya elegido.
\end{note}

Volvamos ahora a la expresión
\[
+\colon\mathbb{N}\times\mathbb{N}\to\mathbb{N}.
\]

La parte
\[
\mathbb{N}\times\mathbb{N}
\]
indica que la operación recibe un par de números naturales. Por eso decimos que la adición es una \emph{operación binaria}: necesita dos números sobre los cuales actuar.

La parte
\[
\to\mathbb{N}
\]
indica que el valor obtenido también es un número natural. En otras palabras, al adicionar dos números naturales, obtenemos nuevamente un número natural.

Esta propiedad recibe el nombre de \emph{clausura}: los números naturales son cerrados respecto de la adición.

\subsection{Operación y operandos}

Supongamos que elegimos dos números naturales cualesquiera, \(a\) y \(b\). Podemos representarlos mediante el par
\[
(a,b).
\]
Al aplicarles la operación de adición, escribimos
\[
a+b.
\]
Los números \(a\) y \(b\) reciben el nombre de \emph{sumandos}. El símbolo \(+\) indica la operación de adición, y la expresión \(a+b\) representa la suma de ambos números.

Si además conocemos otro numeral \(c\) que representa ese mismo número, podemos escribir la igualdad
\[
a+b=c.
\]
Por ejemplo, si \(a=2\) y \(b=3\), entonces
\[
2+3=5.
\]
Aquí, \(2+3\) y \(5\) son dos expresiones que representan el mismo número natural. La primera muestra explícitamente la adición realizada; la segunda expresa su valor mediante un único numeral.

Veamos ahora cómo interpretar estas ideas en un caso concreto.

\section{Práctica de sumas con naturales}

Seguramente esté impresionado con la suma. Sin embargo, de poco le sirve si no sabe aplicarla. Para ello vamos a despejar un poco la mente de las definiciones rigurosas y vamos a sumar algunos números.

\begin{example}
Se le pide realizar la adición entre dos números naturales: \(302\) y \(22\). Como ambos pertenecen a \(\mathbb{N}\), puede aplicar la operación de adición:
\[
\begin{array}{r@{}}
    302 \\
    +22 \\
    \hline
    324 \\
  \end{array}
\]

Conviene detenerse un momento en qué significa matemáticamente esta cuenta. La adición es una operación que recibe dos números naturales y produce otro número natural. Simbólicamente,
\[
+\colon \mathbb{N}\times\mathbb{N}\to\mathbb{N}.
\]
Aquí, el símbolo \(\times\) indica el producto cartesiano\footnote{La multiplicación o producto entre números se estudiará más adelante.}; en este contexto, simplemente expresa que la operación recibe un par de números naturales.

En nuestro caso, ese par es \((302,22)\). Al aplicarle la operación de adición obtenemos
\[
302+22.
\]

Esta expresión ya representa la suma de \(302\) y \(22\). El paso siguiente consiste en \emph{evaluarla}, es decir, encontrar una representación más directa del número que denota:
\[
302+22=324.
\]

La igualdad anterior nos dice que \(302+22\) y \(324\) representan exactamente el mismo número natural. La diferencia está solamente en su forma de escritura: \(302+22\) muestra el número mediante una operación y sus operandos, mientras que \(324\) lo expresa directamente mediante un único numeral decimal.

Así, efectuar la cuenta no crea una nueva cantidad distinta de \(302+22\); permite reconocer y escribir de una manera más simple la cantidad que esa expresión ya representa.
\end{example}

Esta observación permite ver los números de una manera más flexible. Por ejemplo,
\[
302=210+92.
\]
Por lo tanto, en la suma anterior también podríamos sustituir \(302\) por esa expresión y escribir
\[
302+22=(210+92)+22.
\]

De este modo, un mismo número puede estar representado por muchas expresiones diferentes. El numeral \(302\), la expresión \(210+92\) y cualquier otra expresión que tenga el mismo valor representan una misma cantidad.

Esta es también una forma útil de interpretar las cuentas habituales: cuando se presenta una expresión como \(302+22\), la cantidad que representa ya está determinada. Lo que buscamos al efectuar la cuenta es evaluar esa expresión y encontrar una forma más directa de escribir su valor.

A continuación practicaremos algunas adiciones. Aunque se trata de una operación elemental y de uso frecuente, estos ejercicios servirán para familiarizarse con su notación y con la idea de evaluación que acabamos de introducir. Puede utilizar una calculadora para verificar los resultados.

\subsection*{Ejercicios}

En los siguientes ejercicios trabajaremos con distintas formas de representar un mismo número natural. Recuerde que expresiones diferentes pueden representar exactamente la misma cantidad.

\subsubsection*{Evaluar expresiones}

En cada caso, escriba el número natural representado por la expresión dada.

\begin{enumerate}
\item \(3+4\)

\item \(7+0\)

\item \(8+5\)

\item \(23+45\)

\item \(17+17\)

\item \(38+27\)

\item \(24+15+33\)

\item \(186+297\)

\item \(425+368\)

\item \(507+408\)

\item \(2347+1859\)

\item \(3999+2888\)

\item \(47328+96574\)
\end{enumerate}

En cada ejercicio, la expresión presentada ya representa un número natural. Su tarea consiste en encontrar una escritura equivalente mediante un único numeral.

\subsubsection*{Descomponer números}

Ahora realice el proceso contrario. Escriba cada número como la suma de dos números naturales.

\begin{enumerate}
\item \(7\)

\item \(12\)

\item \(25\)

\item \(40\)

\item \(83\)

\item \(125\)

\item \(302\)

\item \(1000\)
\end{enumerate}

No existe una única respuesta correcta. Por ejemplo,
\[
12=7+5,
\]
pero también
\[
12=10+2
\]
y
\[
12=6+6.
\]
Todas estas expresiones representan el mismo número.

Repita ahora algunos de los ejercicios anteriores, pero encuentre tres descomposiciones diferentes para cada número.

\subsubsection*{Completar expresiones equivalentes}

Complete los espacios de manera que ambos lados de la igualdad representen el mismo número.

\begin{enumerate}
\item \(8+4=\square\)

\item \(15=\square+6\)

\item \(20=13+\square\)

\item \(17+8=\square+10\)

\item \(30+12=40+\square\)

\item \(100=60+\square\)

\item \(125=100+\square\)

\item \(302=200+\square\)

\item \(302=210+\square\)

\item \(500+25=400+\square\)
\end{enumerate}

Observe que en algunos ejercicios no se pide simplemente evaluar una suma. Es necesario reconocer que una misma cantidad puede expresarse de distintas maneras.

\subsubsection*{Construir expresiones}

Encuentre una expresión que cumpla con cada condición.

\begin{enumerate}
\item Una suma de dos números naturales que represente \(10\).

\item Una suma de dos números naturales diferentes que represente \(20\).

\item Una suma de dos números naturales iguales que represente \(30\).

\item Tres números naturales cuya suma represente \(25\).

\item Dos expresiones diferentes que representen \(50\).

\item Tres expresiones diferentes que representen \(100\).

\item Una expresión que represente \(302\) sin utilizar el numeral \(302\).

\item Dos expresiones distintas que representen el mismo número natural, sin escribir directamente cuál es ese número.
\end{enumerate}

\subsection{Un algoritmo para evaluar sumas}

Hasta aquí hemos trabajado con sumas que pueden evaluarse mentalmente o mediante descomposiciones sencillas. Sin embargo, cuando los números crecen, resulta útil disponer de un procedimiento sistemático que permita encontrar el valor de una suma sin depender únicamente del cálculo mental.

El método que estudiaremos a continuación es el algoritmo usual de la adición en notación decimal. Este algoritmo no define qué es sumar: la adición ya ha sido presentada como una operación entre números. Su función es ayudarnos a evaluar, de manera ordenada y eficiente, expresiones como
\[
4387+2956.
\]
Para comprender por qué funciona, debemos recordar que la posición de cada cifra dentro de un numeral determina su valor. Por ejemplo,
\[
4387=4000+300+80+7
\]
y
\[
2956=2000+900+50+6.
\]
Por lo tanto,
\[
4387+2956
\]
puede escribirse como
\[
(4000+300+80+7)+(2000+900+50+6).
\]
Si agrupamos las cantidades que ocupan una misma posición decimal, obtenemos
\[
(4000+2000)+(300+900)+(80+50)+(7+6).
\]
Así, el problema puede resolverse trabajando por separado con unidades, decenas, centenas y unidades de millar. El algoritmo usual de la adición aprovecha precisamente esta estructura.

\subsubsection*{Alinear las posiciones}

Para aplicar el algoritmo, los numerales se escriben uno debajo del otro haciendo coincidir las cifras que ocupan la misma posición:
\[
\begin{array}{rrrrr}
 & 4 & 3 & 8 & 7\\
+& 2 & 9 & 5 & 6\\
\hline
\end{array}
\]
Las unidades quedan debajo de las unidades, las decenas debajo de las decenas, las centenas debajo de las centenas, y así sucesivamente.

Esta alineación es fundamental. No estamos colocando las cifras de este modo por una cuestión estética, sino porque queremos sumar entre sí las cantidades correspondientes a una misma posición decimal.

\subsubsection*{Comenzar por las unidades}

Observemos primero la columna de las unidades:
\[
7+6=13.
\]
Sin embargo, \(13\) no puede representarse mediante una única cifra en la posición de las unidades. Debemos descomponerlo:
\[
13=10+3.
\]
Es decir, \(13\) unidades equivalen a \(3\) unidades y \(1\) decena.

Por eso escribimos el \(3\) en la posición de las unidades y trasladamos la decena restante a la columna siguiente:
\[
\begin{array}{rrrrr}
 &   &   & 1 &  \\
 & 4 & 3 & 8 & 7\\
+& 2 & 9 & 5 & 6\\
\hline
 &   &   &   & 3
\end{array}
\]
A este procedimiento se lo conoce habitualmente como ``llevar uno''. Pero ese \(1\) no aparece de manera arbitraria: representa una decena que hemos obtenido al descomponer las \(13\) unidades.

\subsubsection*{Continuar con las decenas}

En la columna de las decenas tenemos ahora
\[
1+8+5=14.
\]
Como cada una de estas cantidades se encuentra en la posición de las decenas, esto representa en realidad
\[
10+80+50=140.
\]
Podemos descomponer las \(14\) decenas como
\[
14\text{ decenas}=1\text{ centena}+4\text{ decenas}.
\]
Por eso escribimos \(4\) en la posición de las decenas y trasladamos \(1\) centena a la columna siguiente:
\[
\begin{array}{rrrrr}
 &   & 1 & 1 &  \\
 & 4 & 3 & 8 & 7\\
+& 2 & 9 & 5 & 6\\
\hline
 &   &   & 4 & 3
\end{array}
\]

\subsubsection*{Continuar con las centenas}

Ahora evaluamos
\[
1+3+9=13
\]
en la columna de las centenas.
Estas \(13\) centenas equivalen a
\[
1\text{ unidad de millar}+3\text{ centenas}.
\]
Escribimos entonces \(3\) en la posición de las centenas y trasladamos \(1\) a la posición de los millares:
\[
\begin{array}{rrrrr}
 & 1 & 1 & 1 &  \\
 & 4 & 3 & 8 & 7\\
+& 2 & 9 & 5 & 6\\
\hline
 &   & 3 & 4 & 3
\end{array}
\]
Finalmente, en la columna de las unidades de millar obtenemos
\[
1+4+2=7.
\]
Así llegamos a
\[
\begin{array}{rrrrr}
 & 1 & 1 & 1 &  \\
 & 4 & 3 & 8 & 7\\
+& 2 & 9 & 5 & 6\\
\hline
 & 7 & 3 & 4 & 3
\end{array}
\]
y, por lo tanto,
\[
4387+2956=7343.
\]

\subsubsection*{¿Qué hace realmente el algoritmo?}

El algoritmo puede parecer una sucesión de reglas sobre cifras, pero en realidad no hace más que organizar descomposiciones como las que ya conocemos.

Cuando una columna produce una cantidad menor que \(10\), esta puede permanecer completamente en esa posición. Por ejemplo,
\[
3+4=7.
\]
Cuando produce una cantidad igual o mayor que \(10\), la descomponemos según el sistema decimal. Por ejemplo,
\[
8+7=15=10+5.
\]
Las \(5\) unidades permanecen en la columna de las unidades, mientras que las \(10\) unidades se reorganizan como \(1\) decena y pasan a la columna siguiente.

Lo mismo ocurre en cualquier posición:
\begin{align*}
  10&\text{ unidades}=1\text{ decena},\\
  10&\text{ decenas}=1\text{ centena},\\
  10&\text{ centenas}=1\text{ unidad de millar},
\end{align*}
y así sucesivamente.

Por esta razón, ``llevar'' una cantidad de una columna a otra no modifica el número representado. Únicamente cambia la forma en que distribuimos esa cantidad entre las distintas posiciones decimales.

\begin{note}
Es importante distinguir nuevamente entre la adición y el algoritmo utilizado para evaluarla.

La expresión
\[
4387+2956
\]
ya representa una suma. El procedimiento vertical que acabamos de estudiar es solamente una herramienta para encontrar una representación más directa de ese mismo número:
\[
4387+2956=7343.
\]

Existen otras formas de llegar al mismo valor. El algoritmo decimal es especialmente útil porque proporciona un procedimiento ordenado que puede aplicarse de la misma manera incluso cuando los números tienen muchas cifras.
\end{note}
